%!TeX program = xelatex
\IfFileExists{RUCReport.sty}{}{\PassOptionsToClass{fontset=fandol}{ctexart}}
\documentclass[12pt,hyperref,a4paper,UTF8]{ctexart}
\IfFileExists{RUCReport.sty}{
    \usepackage{RUCReport}
}{
    % 本地缺少原样式时，沿用原模板的 A4、2.5 cm 边距和标题层级。
    \usepackage[margin=2.5cm]{geometry}
    \ctexset{
        section={format=\centering\zihao{4}\bfseries},
        subsection={format=\zihao{-4}\bfseries}
    }
}
\usepackage{graphicx}
\usepackage{fancyhdr}
\usepackage{hyperref}
\usepackage{needspace}
\usepackage{listings}
\usepackage{xcolor}
% 定义可能使用到的颜色
\definecolor{CPPLight}  {HTML} {686868}
\definecolor{CPPSteel}  {HTML} {888888}
\definecolor{CPPDark}   {HTML} {262626}
\definecolor{CPPBlue}   {HTML} {4172A3}
\definecolor{CPPGreen}  {HTML} {487818}
\definecolor{CPPBrown}  {HTML} {A07040}
\definecolor{CPPRed}    {HTML} {AD4D3A}
\definecolor{CPPViolet} {HTML} {7040A0}
\definecolor{CPPGray}  {HTML} {B8B8B8}
\definecolor{keywordcolor}{rgb}{0.8,0.1,0.5}
\definecolor{webgreen}{rgb}{0,.5,0}
\definecolor{bgcolor}{rgb}{0.92,0.92,0.92}

\lstset{
    breaklines = true,                                   % 自动将长的代码行换行排版
    extendedchars=false,                                 % 解决代码跨页时，章节标题，页眉等汉字不显示的问题
    columns=fixed,
    numbers=left,                                        % 在左侧显示行号
    basicstyle=\zihao{-5}\ttfamily,
    numberstyle=\small,
    frame=none,                                          % 不显示背景边框
    % backgroundcolor=\color[RGB]{245,245,244},            % 设定背景颜色
    keywordstyle=\color[RGB]{40,40,255},                 % 设定关键字颜色
    numberstyle=\footnotesize\color{darkgray},           % 设定行号格式
    commentstyle=\it\color[RGB]{0,96,96},                % 设置代码注释的格式
    stringstyle=\rmfamily\slshape\color[RGB]{128,0,0},   % 设置字符串格式
    showstringspaces=false,                              % 不显示字符串中的空格
    % frame=leftline,topline,rightline, bottomline         %分别对应只在左侧，上方，右侧，下方有竖线
    frame=shadowbox,                                     % 设置阴影
    rulesepcolor=\color{red!20!green!20!blue!20},        % 阴影颜色
    basewidth=0.6em,
}

\lstdefinestyle{CPP}{
    language=c++,                                        % 设置语言
    morekeywords={alignas,continute,friend,register,true,alignof,decltype,goto,
    reinterpret_cast,try,asm,defult,if,return,typedef,auto,delete,inline,short,
    typeid,bool,do,int,signed,typename,break,double,long,sizeof,union,case,
    dynamic_cast,mutable,static,unsigned,catch,else,namespace,static_assert,using,
    char,enum,new,static_cast,virtual,char16_t,char32_t,explict,noexcept,struct,
    void,export,nullptr,switch,volatile,class,extern,operator,template,wchar_t,
    const,false,private,this,while,constexpr,float,protected,thread_local,
    const_cast,for,public,throw,std},
    emph={map,set,multimap,multiset,unordered_map,unordered_set,
    unordered_multiset,unordered_multimap,vector,string,list,deque,
    array,stack,forwared_list,iostream,memory,shared_ptr,unique_ptr,
    random,bitset,ostream,istream,cout,cin,endl,move,default_random_engine,
    uniform_int_distribution,iterator,algorithm,functional,bing,numeric,},
    emphstyle=\color{CPPViolet},
}

\lstdefinestyle{Java}{
    language=[AspectJ]Java,
    keywordstyle=\color{keywordcolor}\bfseries
}

\lstdefinestyle{Python}{
    language=Python,
}

% 封面信息在这里填写。
\newcommand{\KVReportSchool}{}
\newcommand{\KVReportMajor}{}
\newcommand{\KVReportStudentIDs}{}
\newcommand{\KVReportNames}{}
\newcommand{\KVReportDate}{\today}

\fancypagestyle{knowverse}{
    \fancyhf{}
    \fancyhead[C]{\small 知屿 KNOWVERSE\quad 数据库系统概论}
    \fancyfoot[C]{\thepage}
    \renewcommand{\headrulewidth}{0.4pt}
    \renewcommand{\footrulewidth}{0pt}
}
\setlength{\headheight}{15pt}
\providecommand{\reference}{\bibliographystyle{plain}\bibliography{reference}}
\hypersetup{hidelinks,pdftitle={知屿 KNOWVERSE 数据库系统概论课程报告}}

%%-------------------------------正文开始---------------------------%%
\begin{document}

%%-----------------------封面--------------------%%
\hypersetup{pageanchor=false}
\begin{titlepage}
    \thispagestyle{empty}
    \centering
    \vspace*{20mm}
    \IfFileExists{figures/ruc-logo.pdf}{
        \includegraphics[width=.60\textwidth]{figures/ruc-logo.pdf}
    }{
        \includegraphics[width=.60\textwidth]{figures/ruc_logo.eps}
    }\par
    \vspace{43mm}
    {\zihao{1}\bfseries 数据库系统概论课程报告\par}
    \vspace{9mm}
    {\zihao{3}\bfseries 知屿 KNOWVERSE\par}
    \vspace{3mm}
    {\zihao{4}基于学习记录与关系数据库的个人学习助手\par}
    \vspace{25mm}
    {\zihao{-3}\kaishu
        \begin{tabular}{@{}l@{\hspace{3mm}}c@{}}
            学院 & \underline{\makebox[83mm][c]{\KVReportSchool}} \\[4mm]
            专业 & \underline{\makebox[83mm][c]{\KVReportMajor}} \\[4mm]
            学号 & \underline{\makebox[83mm][c]{\KVReportStudentIDs}} \\[4mm]
            姓名 & \underline{\makebox[83mm][c]{\KVReportNames}}
        \end{tabular}\par
    }
    \vspace{10mm}
    {\zihao{4}\KVReportDate\par}
\end{titlepage}

%%--------------------------目录页------------------------%%
\pagestyle{empty}
\tableofcontents
\thispagestyle{empty}
\clearpage

%%------------------------正文页从这里开始-------------------%
\pagenumbering{arabic}
\hypersetup{pageanchor=true}
\pagestyle{knowverse}

\phantomsection
\addcontentsline{toc}{section}{摘要}
\begin{abstract}

\end{abstract}

\section{需求分析}

\subsection{学习场景与功能需求}


\section{概要设计}

\subsection{系统分层与模块分工}


\section{数据库设计}

\subsection{实体、联系与关系模式}

\subsection{候选键、规范化与设计取舍}

\subsection{完整性约束、权限与事务}

\subsection{查询与优化}


\section{核心业务详细设计}

\subsection{个性化答疑与历史检索}

\subsection{资料版本与引用}

\subsection{记忆管理}

\subsection{练习评价与 SQL 沙箱}

\subsection{学习成长与班级协作}

\Needspace{11\baselineskip}

\section{系统实现与交互}

\subsection{主要功能与真实运行示例}


\section{系统测试与实验评测}

\subsection{测试环境与功能验证}

\subsection{数据库测试与查询性能对照}


\section{部署与运行维护}


\section{大模型辅助设计与审查}


\section{总结、局限与后续计划}

%%----------- 参考文献 -------------------%%
% 正文加入 \cite 后，将空标题替换为原模板的 \reference，再运行 BibTeX。
\section*{参考文献}
\phantomsection
\addcontentsline{toc}{section}{参考文献}

\appendix
\section*{附录}
\phantomsection
\addcontentsline{toc}{section}{附录}

\end{document}
